% ===========================================================================
% Figura 1.3: Medición del retardo en gPTP (Peer Delay Mechanism)
% Basada en la Figura 1.3 del documento base
% ===========================================================================
\documentclass[tikz,border=10pt]{standalone}
\usepackage[utf8]{inputenc}
\usepackage[T1]{fontenc}
\usepackage{lmodern}
\usepackage{amsmath}
\usetikzlibrary{shapes.geometric, arrows.meta, positioning, calc}

\begin{document}
\begin{tikzpicture}[
    blk/.style={rectangle, draw, thick, minimum width=2.5cm, minimum height=1.2cm, align=center, font=\small},
    msg/.style={thick, -{Stealth[length=2.5mm,width=2mm]}},
    ts/.style={font=\small, inner sep=2pt}
  ]

  % Nodos
  \node[blk, fill=blue!10] (req) at (0,0) {SOLICITANTE\\(Maestro)};
  \node[blk, fill=green!10] (resp) at (7,0) {RESPONDEDOR\\(Esclavo)};

  % Ejes de tiempo
  \draw[thick, -{Stealth[length=2mm]}] (-0.5,-1) -- (7.5,-1) node[below] {$t$ (solicitante)};
  \draw[thick, -{Stealth[length=2mm]}] (-0.5+7.5,-1) -- (7.5+7.5,-1) node[below] {$t$ (respondedor)};

  % Mensaje 1: Pdelay_Req
  \draw[msg] (req.north) .. controls +(0,4) and +(0,4) .. node[above, font=\small] {\emph{Pdelay\_Req}} (resp.north);
  \draw[fill] (req.north) ++(0,0.5) circle (2pt) node[right, ts] {$t_1$};
  \draw[fill] (resp.north) ++(0,0.5) circle (2pt) node[left, ts] {$t_2$};

  % Mensaje 2: Pdelay_Resp
  \draw[msg, dashed] (resp.south) .. controls +(0,-4) and +(0,-4) .. node[below, font=\small] {\emph{Pdelay\_Resp}} (req.south);
  \draw[fill] (resp.south) ++(0,-0.5) circle (2pt) node[left, ts] {$t_3$};
  \draw[fill] (req.south) ++(0,-0.5) circle (2pt) node[right, ts] {$t_4$};

  % Mensaje 3: Pdelay_Resp_Follow_Up
  \draw[msg, dashed] ($(resp.south)+(0,-1)$) .. controls +(0,-5) and +(0,-5) .. node[below, font=\small] {\emph{PdelayResp}}\node[below, font=\small] {\emph{FollowUp}($t_3$)} ($(req.south)+(0,-1)$);

  % Fórmula
  \node[draw, thick, fill=white, rounded corners, font=\footnotesize] at (3.75,-4.5) {$\displaystyle \text{Pdelay} = \frac{(t_4 - t_1) - r(t_3 - t_2)}{2}$};

  \node[font=\small, text=gray, below, align=center] at (3.75,-5.5) {Fuente: Adaptado de [48].\\$r$ = \textit{rate ratio} = $f_{esclavo} / f_{maestro}$};
\end{tikzpicture}
\end{document}
